% ============================================================
% Chapter 3 -- Preliminary Results
% Numbers come from results/cloud_runs/staged_run_2026-08-03_08-32-45
% and proposal/H2_results.xlsx.
% ============================================================

\paragraph{Setup.}
On PubMedQA \citep{jin2019pubmedqa} with its own 272K-abstract corpus, the dense retriever (\texttt{all-mpnet-base-v2}, 768 dimensions; \citealp{song2020mpnet}) ranked the target first for 93.5\% of the questions, which leaves too little headroom. Following the supervisors' recommendation, we built a harder benchmark. We embedded the full PubMed corpus distributed with MedRAG (${\approx}$23.9M snippets; \citealp{xiong2024medrag}) on a cloud GPU instance. We then interleaved the 100 nearest neighbours of each of 199 questions and of each question's target into a shared pool of 5,035 documents. Every question--document pair (1,000,174 pairs) was scored by Gemini~2.5~Flash-Lite with a rubric on a 1--10 scale, normalised to 0.1--1.0. On this pool B1 ranks the target first for 70.4\% of the 196 analysable questions, so it fails on 58.

\section{The Judge Is a Strong but Coarse Signal (RQ1)}
\label{sec:prelim-judge}

Per query, the judge's score correlates negatively with the distance from the document to the target: the median Spearman correlation is $-0.29$ and the median Pearson correlation is $-0.50$. The correlation is negative and significant for \emph{every} query, and as a detector of the target the judge reaches ROC-AUC 0.997. The judge effectively uses six levels, and the mean distance to the target falls monotonically across them, from 0.89 at score 0.1 to 0.22 at score 1.0. The score is therefore usable as an ordinal ``distance gauge''. Its resolution at the top is poor, however: a median of 18 documents per query share the score 1.0 with the target. Moreover, on 46 of the 58 failed queries the judge also rates B1's wrong first document 1.0, so the current judge rarely flags a retrieval failure.

\section{Geometry of the Query--Document Space (RQ2)}
\label{sec:prelim-geometry}

We drew balls of radius $\dist(q,\GT)/2$ around each query and around its target. The query ball contained no document for all 196 queries, while the target ball was populated for 23\% of them (38\% of the failed ones). Queries are geometric outliers (H1). The centroid of the documents the judge rates 1.0 is 44\% closer to the target than the query (mean distance 0.097 against 0.173), and it is closer for 92\% of the queries, with a larger gain on failed ones. Its ball is smaller than the query's, yet it holds three times more documents. Control experiments showed that random points near the target capture almost nothing, whereas any real document closer to the target than the query captures as many documents as the centroid. Documents are therefore concentrated around the target, and the query lies outside this region. Finally, the query--centroid--target triangle is nearly isosceles with its apex at the query ($\angle q\approx46^{\circ}$): the centroid points towards the \emph{cluster} of relevant documents, not specifically towards the target.

\section{Static and Adaptive Policies (H2)}
\label{sec:prelim-policies}

Table~\ref{tab:policies} compares B1 with policies that mix in the centroid, first as an \emph{oracle} (computed from all judge scores) and then \emph{online} (from the documents judged so far). Sanity checks confirmed that B1 equals the RAG retriever, that the hybrid converges to B1 as $\alpha\to1$, and that a centroid tie-breaker never fires, because cosine distances do not tie.

\begin{table}[htbp]
\centering
\caption{Judge calls until the target is found (mean), on all 196 queries and on the 58 failed ones. Hybrid: $\alpha\,\mathrm{sim}(q,d)+(1-\alpha)\,\mathrm{sim}(C,d)$.}
\label{tab:policies}
\setstretch{1.1}\small
\begin{tabular}{lccc}
\toprule
Policy & Top-1 (\%) & Calls & Calls (failed) \\
\midrule
B1: distance to query (RAG) & 70.4 & \textbf{3.29} & \textbf{8.72} \\
Oracle centroid only & 14.8 & 9.87 & 14.31 \\
Oracle hybrid, $\alpha=0.93$ & \textbf{71.4} & 3.35 & 8.93 \\
Oracle hybrid, $\alpha=0.70$ & 67.9 & 3.71 & 9.86 \\
Top-10, re-ranked by oracle centroid & 17.9 & -- & -- \\
Online hybrid, $\alpha=0.9$ / $0.5$ & 70.4 / 70.4 & 3.35 / 4.10 & 8.93 / 11.47 \\
\bottomrule
\end{tabular}
\end{table}

Even with oracle information, the best hybrid gains one point at top-1 and changes the target's rank in only 16 queries (6 better, 10 worse). Online attraction towards judge-approved documents \emph{increases} the cost, because it leads the search into the dense cluster of relevant-but-non-target documents. Judging B1's top-5 and returning the best (ties broken by distance) improves top-1 only from 73.8\% to 75.4\%, at five calls per query. Three conclusions follow, and they match the action items of our last meeting (12.10.2026): the search should run only when B1 is likely to have failed; it should push \emph{away} from relevant-but-non-target documents rather than towards them; and it should exploit \emph{graded} feedback, which requires a finer-grained judge.
